\section{Preliminary lemmas: character sums, correlation, and monotone bounds}\label{sec:notation}

\paragraph{What this section proves and how to read it.}
This section proves no part of the main argument. It fixes
conventions for finite probability spaces and operators, lists
arithmetic notation, and states five inputs used as black boxes
later: Weil's bound with Gauss and Jacobi evaluations (first-order
$O(p^{-1/2})$ estimates), the spectrum of the square shift on
$\mathbf F_p$, an explicit prime-counting bound, two correlation
inequalities, and a method for extending a bound proved at the
smallest prime of a range to every larger prime (monotone
majorants). Inputs: only standard character-sum and probability
facts, with Peyre's inequality proved in \cref{app:peyre}. Outputs:
the lemmas cited by Parts~IV and~V for counting, correlation after
fixing labels, and prime summation. A first reading should note the
conventions in \cref{subsec:notation:conventions} and the statements
only; proofs can be skipped. The list in
\cref{subsec:notation:symbols} is an \emph{index of symbols} with
pointers to defining sections, not a set of definitions; project
terms are defined in \cref{subsec:intro:glossary} and at their formal
sites. In particular ``price'' is defined in \cref{def:price}, not
here.

\subsection{Finite probability spaces, operators, and correlation}\label{subsec:notation:conventions}

All probability spaces in the paper are finite. This subsection
fixes expectation, norms, operator order and correlation once, so
that later sections can use them without repeating conventions.
Asymptotic notation is the standard one, made precise as follows:
$X\ll Y$, $X=O(Y)$ and $X\lesssim Y$ all mean $|X|\le CY$ for an
absolute constant $C$, independent of $N$, of the scale $\Lsc$, of
$A$ and of $p$, unless a dependence is stated in the same
statement or in the standing conventions of its section; $X\asymp
Y$ means both $X\ll Y$ and $Y\ll X$; and $o(1)$ denotes a quantity
tending to $0$ as $\Lsc\to\infty$ with all other parameters fixed,
unless another limit is named. In statements at a least
counterexample, $o(1)$ is this limit along counterexamples with
$\Lsc\to\infty$.

\begin{definition}[Finite probability spaces and operator order]\label{def:prelim:operators}
For a finite probability space $(X,\mu)$,
$\mathbb E_\mu f=\sum_{x\in X}\mu(x)f(x)$ and
$\|f\|_{L^r(\mu)}=(\mathbb E_\mu|f|^r)^{1/r}$ for $1\le r<\infty$.
The inner product is $\langle f,g\rangle_\mu=\mathbb E_\mu\overline f g$;
$L^2_0(\mu)$ is the subspace of mean-zero functions.
All unqualified averages on finite sets use uniform probability.
A measure $\nu$ has density $h$ with respect to $\mu$ if
$\nu(x)=h(x)\mu(x)$ at each point; its mass is $\nu(X)$. The notation
$\delta_x$ denotes the probability measure concentrated at $x$.
For a self-adjoint operator $T$, the notation $T\succeq0$ means
$\langle f,Tf\rangle\ge0$ for every $f$. We write $S\preceq T$
if $T-S\succeq0$; this is the \emph{Loewner order}.
The operator norm is $\|T\|=\sup_{\|f\|_2=1}\|Tf\|_2$.
For a numerical matrix, the operator norm uses the ordinary Euclidean
norms on its coordinate spaces: matrix multiplication is a sum, with
no probability normalization. In particular, this convention applies
to the matrix of correlation bounds in \cref{thm:peyre-finite}.
A \emph{Markov operator} preserves nonnegative functions and the
constant function $1$.
\end{definition}

\paragraph{Kernels, projections and centred norms.}
Probability spaces are finite and, after the points of mass zero are
discarded, have full support. For such a space $(\Omega,\lambda)$,
$\Jproj f=\ip{\one}{f}_\lambda\one=(\E_\lambda f)\one$ and
$P_0=I-\Jproj$ (in Part~IV also written $Q$; this $Q$ is not the
conductor cutoff $Q$ of \cref{sec:lifts}) are the
orthogonal projections onto the constants and onto
$L^2_0(\lambda)$. For finite probability spaces $(X,\mu)$, $(Y,\nu)$ and a
kernel $K\colon X\times Y\to\C$ let $T_Kf(x)=\int K(x,y)f(y)\,d\nu(y)$, and let
$\opnorm K$ be the norm of $T_K$ from $L^2(\nu)$ to $L^2(\mu)$. Both spaces carry
the uniform probability unless another measure is indicated; for the uniform
probabilities, $\opnorm K$ is the operator norm of the entry matrix of $K$
(Euclidean norms) divided by $\sqrt{|X||Y|}$. If
$X=Y$, $\mu=\nu$ and $K(y,x)=\overline{K(x,y)}$, then $T_K$ is self-adjoint. The
all-ones kernel is $\Jall$, so that $T_{\Jall}=\Jproj$ on a uniform space; as a
matrix acting on vectors indexed by $n$ points, $\Jall v=(\sum_yv(y))\one=n\Jproj v$.

Let $T$ be a Markov operator from functions on $Y$ to functions on $X$ with
$T\one=\one$ and $T^*\one=\one$, the adjoint taken with respect to $\mu$ and
$\nu$. Equivalently, $Tf(x)=\E_\pi[f(y)\mid x]$ for the coupling
$\pi(x,y)=\mu(x)T(x,y)$ of $\mu$ and $\nu$, $T(x,y)$ being the transition
probability, so that $Tf(x)=\sum_{y\in Y}T(x,y)f(y)$; we call $T$ the operator of the coupling. As
$\ip{\one}{Tf}_\mu=\ip{T^*\one}{f}_\nu$, $T$ maps $L^2_0(\nu)$ into $L^2_0(\mu)$,
and its \emph{centred norm} is $\|T|_{L^2_0(\nu)}\|=\|T-\Jproj\|$, with
$\Jproj f=(\E_\nu f)\one$. By \cref{def:prelim:maxcorr} the centred norm of the
operator of a coupling is the maximal correlation of its two coordinates. For a
cut $c$ of a slot array, $\kappa\succeq_c\kappa'$ means that the cut matrix of
$\kappa-\kappa'$ is symmetric and positive semidefinite (\cref{def:cc:order}).

\begin{definition}[Mean, fluctuation and compression]\label{def:split:notation}
On a finite probability space $(\Omega,\lambda)$ with full support let
$\Jproj f=(\E_\lambda f)\one$ and $P_0=I-\Jproj$. For a real function $U$ put
$u=\E_\lambda U$, $U^\circ=P_0U$ and $\sigma_U=\norm{U^\circ}_2$. Thus
$\norm U_2^2=u^2+\sigma_U^2$. For an operator $T$ on $L^2(\lambda)$ with
$\ip\one{T\one}=1$, its two \emph{marginal errors} are $c=T\one-\one$ and
$b=T^*\one-\one$; its \emph{mean-zero compression} is $C=P_0TP_0$. Relative to
the constants and their orthogonal complement,
\[
 T=\begin{pmatrix}1&b^*\\c&C\end{pmatrix},
\]
where $b^*$ is the functional $g\mapsto\ip bg$ on $L^2_0(\lambda)$.
\end{definition}

Both marginal errors are mean-zero, since
$\ip\one c=\ip\one{T\one}-1=0=\ip{T^*\one}\one-1=\ip b\one$, and for
$f=\alpha\one+f_0$ with $f_0\in L^2_0(\lambda)$,
$Tf=(\alpha+\ip b{f_0})\one+(\alpha c+Cf_0)$, which is the block form. The
operator of a coupling of $\lambda$ with itself has $b=c=0$ and centred norm
$\|C\|$. The completion inequalities of \cref{lem:five-type} are written in
this notation.

\begin{definition}[Maximal and conditional maximal correlation]\label{def:prelim:maxcorr}
For random variables $X,Y$ on a finite probability space, the \emph{maximal
correlation} is
\[
  \rho(X;Y)=\sup\big\{|\E f(X)g(Y)|:\ \E f(X)=\E g(Y)=0,\ \E f(X)^2=\E g(Y)^2=1\big\},
\]
with $\rho=0$ if either side has no non-constant functions. The intersection
of the spaces $L^2(X)$ and $L^2(Y)$ of functions of $X$ and of $Y$, inside
$L^2$ of the underlying space, contains the constants. If it consists of the constants only, then
$\rho(X;Y)$ is the Friedrichs cosine of these two subspaces
(the supremum of $|\langle f,g\rangle|$ over unit vectors $f$ in
one subspace and $g$ in the other, after the constants are
removed); if it is larger,
then $\rho(X;Y)=1$, while the Friedrichs cosine may be smaller. In all cases
$\rho(X;Y)$ is the norm of the conditional expectation operator
$L^2_0(X)\to L^2_0(Y)$. For a family $(X_i)_{i\in I}$ and $S\subseteq I$, the
\emph{conditional} maximal correlation of $X_i,X_j$ given $X_S=s$ is the
maximal correlation computed under $\Pr(\cdot\mid X_S=s)$, for every value
$s$ of positive probability.
\end{definition}

\begin{definition}[Friedrichs angle]\label{def:prelim:angle}
For closed subspaces $V,W$ of a finite-dimensional Hilbert space, the
\emph{cosine of the Friedrichs angle} is
$\cos\angle(V,W)=\sup\{|\ip vw|:\ v\in V,\ w\in W,\ \norm v=\norm w=1,\
v,w\perp V\cap W\}$. It is $0$ if one space contains the other.
\end{definition}

\subsection{Arithmetic notation, square classes, and index of symbols}\label{subsec:notation:symbols}

\paragraph{Arithmetic notation.}
Throughout, $p$ denotes a prime, $q=1/p$ and $r=p^{-1/2}$; for odd $p$,
$\chi=\chi_p$ is the quadratic character. We write $[N]=\{1,\dots,N\}$ and
$\ee t=e^{2\pi it}$. All logarithms are natural, an empty product is $1$, and
a positive integer is \emph{squarefree} if no square of a prime divides it.
For a positive integer $d$ we write $v_p(d)$ for the exponent of $p$ in $d$,
$\varphi(d)$ for the number of invertible residues modulo $d$, $\omega(d)$ for
the number of its distinct prime divisors (always with its argument) and
$\rad(d)$ for their product; $\pi(x)$ is the number of primes up to $x$. For
$\xi\in\Q/\Z$, $q(\xi)$ is its reduced denominator, the \emph{conductor}, with
$q(0)=1$; for a set $U$ of primes, $q_U=\prod_{p\in U}p$. For an odd prime $p$
we put
\[
  \Qsq=\{u^2:u\in\Fp^\times\},\qquad \Ap(t)=\one_{\Qsq}(t),\qquad
  \Om=\{(x,y)\in\Fp^2:\ y-x\in\Qsq\};
\]
$\chi_p(t)$ is $1$ on nonzero squares, $-1$ on non-squares and $0$ at $t=0$,
and $G_p(k)=\sum_{u\in\Fp}\ee{ku^2/p}$. The space $L^2(\Fp)$ carries the
uniform probability, $\ip fg=\E_{a\in\Fp}\overline{f(a)}g(a)$, and
$e_k(a)=\ee{ka/p}$, $k\in\Fp$, is an orthonormal basis; the $e_k$ with
$k\ne0$ span $L^2_0(\Fp)$. For a law $\nu$ on a finite set and functions $F$ and $G$ on that set,
$\norm F_\nu^2=\E_\nu|F|^2$ and $\ip FG_\nu=\E_\nu\overline FG$; the index is
dropped when the law is clear.

\begin{definition}[Strict square classes]\label{def:strict-square-class}
Let $D=8d$ with $d$ odd, not necessarily squarefree. A residue $r\bmod D$ is a
\emph{strict square class} if $r\equiv1\pmod8$ and $r$ is a nonzero square
modulo every prime dividing $d$.
\end{definition}

By Hensel's lemma the second condition is equivalent to $r$ being the square
of a unit modulo every odd prime power dividing $D$: if $p$ is odd, $p\nmid r$ and $c_j^2\equiv
r\pmod{p^j}$, then $c_{j+1}=c_j+tp^j$ with $2c_jt\equiv(r-c_j^2)/p^j\pmod p$
satisfies $c_{j+1}^2\equiv r\pmod{p^{j+1}}$, and $t$ exists because
$p\nmid2c_j$. Every unit modulo $8$ has square $1$, so by the Chinese remainder
theorem the strict square classes modulo $D$ are exactly the squares of units
modulo $D$; in particular a strict square class is a unit modulo $D$.

The external index of \cref{def:hyb:coupled} also uses nonunit squares at a fixed prime-power
modulus. We use the following more general class only when explicitly stated.
\begin{definition}[Nonvanishing square classes]\label{def:nonvanishing-square-class}
Let $D=2^{h_2}\prod_{p\text{ odd}}p^{h_p}$ with $h_2\ge3$.
A residue $r\bmod D$ is a nonvanishing square class if, at each odd
$p\mid D$, it has valuation $2j_p<h_p$ and square unit part, and at $2$
it has valuation $2j_2$ with $h_2-2j_2\ge3$ and unit part congruent
to $1$ modulo $8$. The valuations and unit parts are properties of the
reductions of $r$ modulo $p^{h_p}$ for odd $p$ and modulo $2^{h_2}$ at $2$. In particular zero modulo a
prime-power coordinate is excluded.
\end{definition}
These conditions imply the existence of a square root modulo each full
prime-power coordinate. At odd primes this is Hensel lifting after
division by $p^{2j_p}$. At $2$, begin with an odd root modulo $8$: if
$b^2\equiv u\pmod{2^k}$ with $k\ge3$, one of $b$ and $b+2^{k-1}$
squares to $u$ modulo $2^{k+1}$, since their squares differ by
$2^kb$ modulo $2^{k+1}$. This lifts every unit $u\equiv1\pmod8$.
Strict square classes with $v_2(D)=3$ are a special case.

\paragraph{The least square step.}
For a positive integer $d$ put
\begin{equation}\label{eq:sq}
 \sq(d)=\prod_{p\mid d}p^{2\lceil v_p(d)/2\rceil},
\end{equation}
the least positive square divisible by $d$: a square $\prod_pp^{2k_p}$ is
divisible by $d$ exactly when $2k_p\ge v_p(d)$ for every $p$. Thus
$d\mid\sq(d)\mid d^2$, $\sq(de)=\sq(d)\sq(e)$ for coprime $d,e$, and
$\sq(p)=p^2$, $\sq(8)=16$, $\sq(24)=144$.

\paragraph{Conventions on numbers.}
$o(1)$ is as defined in \cref{subsec:notation:conventions}; in
the analytic statements of Part~II the convergence as
$\Lsc\to\infty$ is uniform in the set $A$. Numbers in the text
are the numbers certified in \cref{app:certificates}, written as in the
certificate (outward-rounded decimals or exact rationals); a decimal that is
not certified is marked ``about''. A \emph{price} is $\|\zeta\|_{\rm
var}q_{\Blk}^{2a}$ for a signed block measure $\zeta$ (\cref{def:price}); an
\emph{operator price} is a weight times an operator norm relative to the
actual half marginals times $q_{\Blk}^{2a}$; an \emph{upper price} is any
number at least the price, in practice a certified upper bound. Every coefficient that is
\emph{evaluated at an endpoint} is an admissible majorant in the sense of
\cref{lem:majorant}, written with its powers of $r$ cancelled first
(\cref{rem:majorant-use}).

\paragraph{Symbols.}
\Cref{tab:notation} lists the principal symbols and where they are defined.
Its rules concern named quantities, written with the bare letter: $t$ denotes
only the repricing factor, $\tau$ and $\sigma$ only the analytic-stage
parameters, $\iota$ only the reserved cut and $J$ only the matching term; and
$\delta_{p,j}$ at a prime-power index coordinate denotes its depth tail,
while $\delta_p$, $\delta_j$ denote inclusion probabilities in the
whole-coordinate special case. A letter used as a bound
variable inside one statement or proof is local to it, for instance a number
$t$ of roots or the index of $e_t$, a type $\tau$, a permutation $\sigma$, an
index set $J$, the matrix $\mathcal E$ of correlation bounds in
\cref{thm:peyre-finite}, or the point mass $\delta_x$. In the same local
sense $m$ or $m_0$ may denote a probability measure on a set of labels
(\cref{sec:labellings,sec:middle-records,sec:envelopes}), $m$ a multiset of
middle pieces (\cref{lem:piece-caps} and
\cref{sec:variance,sec:envelopes,sec:matching-mean}), and $t$ may index
the records of a block (\cref{sec:comparison-theorem,sec:comparison-lemmas}
and Part~V); the mean credit $m$ occurs only in the closing formulas, and the
repricing factor $t$ only in formulas that carry a price certified at the
anchor $a_0$ to the exponent $a$: the repricing statements
(\cref{lem:record:repricing,cor:first-stage-repricing}), the totals and
credits of Part~V formed from such prices (for instance the middle prices,
the variance credit and the class costs of the envelopes,
\cref{def:closing-quantities,cor:envelope-costs}) and the closing formulas;
these local uses do not occur there. Two rules are followed in every section: the letter $B$ alone is the energy
factor of the reserved Gram bracket, and a block of primes is written $\Blk$
wherever both can occur (Part~IV, which never meets the energy factor, may
write a block as $B$); and the layer coordinate of a slot is $\ell$, so that
$t$ is free for the repricing factor. Several objects of the local laws carry
sans-serif letters: the Paley channels $\mathsf A$--$\mathsf D$ of the fold
record, the scores $\Hs j$ of the positive roots, the signed-law parameter
$\tlaw$ and the root-tilt parameter $\ulaw$; neither $\tlaw$ nor $\ulaw$ is
ever the repricing factor $t$. The reserved cost is $\mathcal E$
(calligraphic, without subscript), and the query family is $\mathcal E_6$
(always with its subscript).

\begingroup
\small
\setlength{\LTpre}{6pt}\setlength{\LTpost}{6pt}
\begin{longtable}{@{}>{\raggedright\arraybackslash}p{0.22\textwidth}%
>{\raggedright\arraybackslash}p{0.74\textwidth}@{}}
\caption{Principal symbols. Numerical illustrations are at the baseline
parameters unless the selected parameters are specified.}\label{tab:notation}\\
\toprule
symbol & meaning (and where it is defined)\\
\midrule
\endfirsthead
\toprule
symbol & meaning (and where it is defined)\\
\midrule
\endhead
\bottomrule
\endfoot
$N$, $[N]$, $s(N)$ & the scale of the main theorem; $\{1,\dots,N\}$; the
 largest size of a subset of $[N]$ with no nonzero square difference\\
$A$, $\alpha$ & a square-difference-free set $A\subseteq[N]$ and its density
 $\alpha=|A|/N$; the letter $A$ alone is also the skew total (row $S$, $A$,
 $T$)\\
$\epsilon$ & the exponent of the main theorem, $s(N)\le CN^{1-\epsilon}$,
 $\eps=a/(72-\theta)=\FourRoutesExponent$; the baseline has
 $\eps=a/71=\PushExponent$\\
$a$, $a_0$, $a_{\rm u}$ & the price exponent of the functional,
 baseline $a=\PushA$, selected $a=\FourRoutesA$; the anchor $a_0=\PushAnchor$, at
 which the middle prices are certified (\cref{subsec:closing-quantities}); the
 upper ledger has its own anchor $a_{\rm u}=883/10^6$
 (\cref{prop:record:upper})\\
$\theta$ & the density exponent of the induction: baseline $1$, selected
 $\FourRoutesDensity$ (also denoted $\vartheta$ in the closing);
 chamber reflections carry their cut indices, and $\theta(x;4,r)$
 always carries its arguments\\
$\Lsc$, $\Kint$ & the interval scale, $N=\Kint\Lsc$; the number of intervals
 of the index: baseline $\Kint=\PushK$, selected $\Kint=\FourRoutesK$\\
$L$, $L_{\rm loc}$, $L_{\rm band}$, $L_{\rm u}$ & the actual block budget
 $L=570$ of the middle range (\cref{def:record:exact-blocks}), used for the
 repricing and for the caps in the restorations of \cref{lem:record:restoration}; the local cap $L_{\rm loc}=585$, at which
 the local laws, the block densities, the row constants, the spectral caps and
 $k_{\rm D}$, $k_{\rm new}$ are evaluated at every piece start
 (\cref{lem:record:local-domain}); the budget $L_{\rm band}$ of an upper band,
 an integer between $550$ and $886$ at which the local data of the band are
 evaluated (\cref{def:record:upper-bands}); and the block length
 $L_{\rm u}=50000/27$ of the local data of the tail datum, larger than every
 band budget (\cref{def:record:parameters})\\
$G_N$, $e_p$ & the label group $\prod_p\Z/p^{e_p}\Z$ over the set
 $\mathcal P_N$ of primes of the scale $N$, with $e_p=2\mcl$ at the
 primes $p\equiv1\ (4)$ below $\Pfour$ and $e_p=1$ otherwise
 (\cref{def:record:functional})\\
$V$, $d$ & the slot array $(\Z/6\Z)^2\times\{0,1\}$ with slots
 $(u_0,u_1,\ell)$, $d=|V|=72$; the layer coordinate is $\ell$
 (\cref{def:slot-array})\\
$\rot$ & the rotation of $\Z/6\Z$: $\rot j=j+2$ for even $j$ and $\rot j=j-2$
 for odd $j$\\
$\vartheta_{c,a}$, $H_a$ & the chamber reflection $\vartheta_a(j)=2a-1-j$ in the
 coordinate $c\in\{0,1\}$, with plus half $H_a=\{a,a+1,a+2\}$\\
$\iota$ & the reserved cut $(u_0,u_1,\ell)\mapsto(u_0,u_1,1-\ell)$, with plus half
 $\{\ell=0\}$; $\iota$ is used for nothing else, and the class indicator is
 written $\one_{p\equiv1\,(4)}$\\
$\mathcal E_6$, $e=(u_e,v_e)$ & the six queries $u_e=(0,j,0)$,
 $v_e=(\rot0,\rot j,1)$, $j\in\Z/6\Z$ (\cref{def:slot-array})\\
$p$, $q$, $r$; $\chi$, $\Qsq$, $\Om$ & a prime, $q=1/p$, $r=p^{-1/2}$; the
 quadratic character, the nonzero squares, the strict-square pairs\\
$G_p(k)$, $G_\chi$, $\eps_G$, $\Jac$ & Gauss sums, $G_\chi^2=\chi(-1)p$,
 $\eps_G=G_\chi/\sqrt p$ (\cref{thm:gauss}); the Jacobi sum
 (\cref{thm:jacobi})\\
$\delta_{p,j}$; $\delta_p$, $\delta_j$ & the conductor-depth tail
 $\Pr(J_{e,p}\ge j)$ of \cref{def:hyb:coupled}, with the selected
 $m=64$ law of \cref{lem:imp:padic-envelope}; the latter two symbols
 are inclusion probabilities in the whole-coordinate special case\\
$\omega_e=(b_e,s_e)$ & the index of the query $e$: an interval and a residue
 modulo $\Mzero$ (\cref{def:hyb:coupled}); count errors are written $e_t$,
 $w_\Delta$, never $\omega$\\
$\Mzero$, $D_\Jc$, $\sq(d)$ & the external modulus
 $2^{129}\prod_{p<Q_3,\,p\equiv3\,(4)}p^{128}$ at $m=64$;
 the revealed modulus $\prod_p p^{J_{e,p}}$ for one query;
 the least square divisible by $d$, \eqref{eq:sq}\\
$\Jc$, $R_I$ & the independent depth vector, with every interval revealed
 (\cref{def:hyb:coupled}); the conductor moment
 $\E_\Jc(\Kint\sq(D_\Jc))^{a/36}=M_{\rm pad}(\Kint)$,
 bounded above by $1.247140$ (\cref{cert:index-moment})\\
$Q_3$, $\Pthree$, $\Pfour$, $\Pfive$ & $19951586$, $167\cdot2^{40}$, $2^{512}$,
 $2^{65536}$: the primes $p\equiv3\ (4)$ of $[Q_3,\Pthree)$ form the middle
 range and those of $[\Pthree,\Pfour)$ the upper range; the large primes are those from $\Pfour$ on, one atom each; mixing in
 the sense of \ref{it:record:R4} is asserted only from $\Pfive$ on
 (\cref{def:record:parameters})\\
$\mcl$, $m$ & the number of digit pairs of the square-clique prefix,
 $\mcl=8192$ (\cref{def:record:parameters}); the mean credit in $\Theta$
 (\cref{def:cp7})\\
$B$ & only the energy factor of the reserved Gram bracket,
 $B=e^{U_0}+R_{\rm err}$ (\cref{def:cp3}); its certified upper end at the
 baseline parameters is $2.332628665480$\\
$\Blk$, $q_\Blk$, $b_\Blk$ & a block (atom) of primes; the product and the
 least of its primes\\
$\Nbr$, $\Nbrz$ & the reserved Gram bracket measures (\cref{def:cp3}), never
 the scale $N$\\
$S$, $A$, $T$ & the symmetric, skew and matching price totals
 (\cref{thm:comparison})\\
$P$, $V$, $D$, $U_0$, $R_{\rm err}$, $L_0$ & the first-stage and bracket
 totals, with the tag totals $d_j$, $v_j$; $L_0=e^P\tagser(d,v)$ and
 $R_{\rm err}=e^{U_0}\tagser(d,v)$ (\cref{def:cp2,def:cp3}); for two tags with
 equal totals the series $\tagser$ has the finite majorant $\tagser_n^+$
 (\cref{cor:two-cut-series})\\
$\mathcal E$, $J$, $W$, $\Theta$ & the reserved cost, the matching term, the
 wrap and the closing number,
 $\Theta=\frac{63}{64}(L_0+\mathcal E)+J+W+\frac{1-m}{64}$ for six queries
 in the baseline (\cref{thm:comparison}); the selected joint formula is
 \eqref{eq:final-joint-closing}; $\Theta(z)$ is its inflated form\\
$D$, $D_{\rm var}$ (closing) & the total discount in
 $\mathcal E=B(e^S\cosh A-1)-D$, the sum of the variance credit $D_{\rm var}$
 and the savings of the envelope classes (\cref{def:closing-quantities}),
 never in one formula with the first-stage total $D$; the first-stage sharp
 total enters the closing only through its two equal tag totals $d_0=d_1$,
 each one half of it (\cref{def:closing-quantities})\\
$\mathsf A$, $\mathsf B$, $\mathsf C$, $\mathsf D$ & the four Paley channels
 $\chi(y_2-x_0)$, $\chi(x_2-y_0)$, $\chi(y_0-x_1)$, $\chi(y_1-x_0)$ of the $F$
 record, sans-serif so as not to collide with the prices\\
$t$ & only the repricing factor $e^{2(a-a_0)L}$, about $1.0000000616512$
 ($a_0=\PushAnchor$, $L=570$; \cref{def:closing-quantities})\\
$\beta$, $\eta$, $\gamma$, $\tau$, $t_{\rm an}$, $\sigma$ & the analytic-stage
 parameters $1/20$, $17/100$, $255/1024$, $1/201$, $1/60$, $1/825$
 (\cref{def:record:analytic-stage})\\
$\tside$, $\tau^{\rm sp}$, $t_Y$, $\Sigma_i$ & the direct-Gram side weight
 $31/32$; the free parameter of the spectral test; the mean (Young) allowance
 of a row of the upper completion, free in the joint rows of the bands
 (\cref{lem:five-type}(f)) and $1/50$ in the Schur rows of the tail datum
 (\cref{lem:five-type}(c)); the block-start sums, certified lower bounds for $\sum_{\Blk}b_\Blk^{-1}$ over the blocks of piece $i$ (\cref{lem:record:lower-start-sums})\\
$\Jall$, $\Jproj$ & the all-ones kernel; the projection onto constants; $J$
 alone is the matching term\\
$E_0$, $E_1$, $F_0$, $F_1$, $F_2$, $G$, $H$, $C$ & $E_0$ and $E_1$ matching records, $F_0$, $F_1$ and $F_2$ fold records,
 $G$ merge records, $H$ reserved half-diagonal records and $C$ constant records of the local laws (\cref{def:patterns,def:record:frame})\\
$e_t(p,n)$, $w_\Delta$, $e^\pm_t$, $a_t$, $b_t$, $h_t$ & the coarse and graded
 insertion-count errors and their coefficients
 (\cref{eq:graphs:et,eq:graphs:graded}, \cref{lem:extension})\\
$A_F$, $A^{\rm arm}_F$, $A^{\rm lead}_F$, $A^{\rm rem}_F$, $S_F$ (and $E$);
 $S_{\rm oth}$, $A_{\rm oth}$ & the full skew price, the raw-arm price (coefficient $w_{\rm D}r^2$ per fold
 record, before the leading and remainder split), and the leading and remainder prices and the symmetric price of the fold (matching) records, with
 $A_F=A^{\rm lead}_F+A^{\rm rem}_F$, while the full skew price $A_E$ of the
 matching records is bounded directly; the other middle symmetric and skew
 prices (\cref{lem:middle-comparison-prices}, \eqref{eq:middle-piece-prices})\\
$\kappa_{FF}$ & the cap of the envelope of the pairs of fold records, the
 fraction of their variance gain that is credited (\cref{lem:class-credit})\\
$\tlaw$ & the signed-law parameter $\tlaw=\FurtherT$ of the fold records
 (\cref{def:record:mixture})\\
$\mathfrak v_\nu$, $n_\nu$ & the variation factor of the comparison law and
 the number of its product components (\cref{def:cp4})\\
$\ulaw$; $\Hs0$, $\Hs1$, $\Hs2$, $\Hbar$ & the root-tilt parameter, with
 its seven values for the seven tilted pieces in \cref{def:record:parameters};
 the three opposite-base-pair scores $\Hs0$, $\Hs1$, $\Hs2$ and their mean
 $\Hbar$. The last of the seven tilted pieces uses the signed coherent trace
 sum $H_{\mathcal D_{39}}$ and the positive floor $1/100$
 (\cref{def:positive-root})\\
$\tagser(d,v)$, $d_j$, $v_j$ & the tagged first-stage series and the sharp
 and variation totals of the tag $j$ (\cref{def:cp2})\\
$\Uinv$ & the assigned price of the invariant tags $E_{\rm add}+H$ of the
 outside law $\Nbrz$ (\cref{lem:oct24:invariant-background})\\
$C_{n,m}$, $\mathcal R_{n,m}$, $D_{n,m}$ & the leading class cost per raw-arm
 monomial, the remainder polynomial and the saving of the class with $n$ fold
 and $m$ matching skew records (\cref{cor:envelope-costs})\\
$M_{n,m}$, $M^{\rm inv}_{n,m}$, $\delta_{\rm s}$ & the collected coupling mass
 and the invariant orbit mass of a class; the price saving of one folded tag,
 sharp minus symmetric (\cref{lem:orbit-coupon})\\
$m_F$, $m_Q$ & the fold and nonempty-$Q$ coupons of the mean credit $m$
 (\cref{prop:record:mean-credit})\\
$\Phat$ & the normalized Paley operator $P/r$ at a class-three prime;
 $\Phat^*=-\Phat$, and $\Phat^2=-I$ on $L^2_0$ (\cref{lem:square-shift}(c))\\
$s_\tlaw$, $d_\tlaw$, $\ell_\tlaw$, $e_{\rm sg}$ & the ideal and actual
 coefficients of the signed direct Gram, the leading cap over $r$ and the
 signed remainder constant (\cref{lem:record:direct-gram,prop:signed-bounds})\\
$k_{\rm D}$, $e_o$; $k_{\rm new}$, $s_{\rm fold}$; $\bar H_X$ & the negative
 cap and the centred remainder coefficient of the reference fold record
 (\cref{lem:old-f-law}); the sharp and symmetric coefficients of the one-stage
 folded repair (\cref{lem:record:one-stage-folded-e}); the block density
 factor $(1-n_Xe_X)^{-1}$ of a piece starting at $X$\\
$\delta^+_{s,i}$ & the increment of the side $s\in\{L,R\}$ of the piece $i$
 over its auxiliary side, $\max(0,\bar H_{X_i}d_\tlaw(X_i)-d^{\rm aux}_s(X_i))$
 rounded up (\cref{prop:signed-sides}); not to be confused with $\delta_{\rm s}$\\
$g$, $c^+$, $B_{\rm top}$ & the prime caps of a cut,
 $-gp^{-1/2}I\preceq\Op\preceq(1+c^+/p)I$ for its normalized cut operator $\Op$ at $p$;
 the positive constant is written $c^+$ for the cuts of the middle range and
 $B_{\rm top}$ for the upper range, where the caps come from Route~B
 (\cref{lem:layer12-caps,thm:record:direct-star,lem:spec:tensorcomp})\\
piece, band & a range of the middle partition (\cref{def:record:exact-blocks});
 a range of the upper partition (\cref{def:record:upper-bands})\\
\end{longtable}
\endgroup

\subsection{Character sums}\label{subsec:prelim:characters}

Let $p$ be an odd prime. A \emph{multiplicative character} is a homomorphism
$\psi\colon\Fp^\times\to\C^\times$; we extend every multiplicative character,
the trivial one $\psi_0$ included, by $\psi(0)=0$, and only nontrivial
characters will enter the sums below. The characters form a cyclic group of
order $p-1$ (the dual of the cyclic group $\Fp^\times$); $\chi_p$ is its
unique element of order $2$. Hence for every $m\mid p-1$ there are characters
of order exactly $m$, and if $4\mid p-1$ every character $\psi$ of order $4$
satisfies $\psi^2=\chi_p$. Since $|\psi(t)|=1$ for $t\ne0$,
$\psi(t^{-1})=\overline{\psi(t)}=\bar\psi(t)$. The element $-1$ is the
element of order two of the cyclic group $\Fp^\times$, so it is a square
exactly when $4\mid p-1$: $\chi(-1)=1$ if $p\equiv1\pmod4$ and $\chi(-1)=-1$
if $p\equiv3\pmod4$. For $s\in\Fp$ the \emph{Gauss sum} of $\psi$ is
\[
  g_s(\psi)=\sum_{t\in\Fp}\psi(t)\,\ee{st/p},\qquad g(\psi)=g_1(\psi),\qquad
  G_\chi=g(\chi_p)=\sum_{t\in\Fp}\chi_p(t)\,\ee{t/p},
\]
and for multiplicative characters $\psi_1,\psi_2$ the \emph{Jacobi sum} is
$\Jac(\psi_1,\psi_2)=\sum_{t\in\Fp}\psi_1(t)\psi_2(1-t)$.

\begin{theorem}[Weil's bound]\label{thm:weil}
Let $p$ be an odd prime and $a_1,\dots,a_t\in\Fp$ distinct, $t\ge1$.
\begin{enumerate}[label=\textup{(\alph*)}]
\item $\big|\sum_{x\in\Fp}\chi_p\big((x-a_1)\cdots(x-a_t)\big)\big|\le(t-1)\sqrt p$.
\item If $t\ge2$ is even, then
 \[
  \Bigl|\sum_{x\in\Fp}\chi_p\bigl((x-a_1)\cdots(x-a_t)\bigr)+1\Bigr|\le(t-2)\sqrt p ,
 \]
 so the sum has modulus at most $(t-2)\sqrt p+1$.
\end{enumerate}
\end{theorem}

Part (a) is Weil's theorem; see \cite[Theorem~11.23]{IwaniecKowalski},
\cite[Ch.~II]{Schmidt76}. Part (b) follows from (a).

\begin{proof}[Proof of \textup{(b)}]
Put $f(x)=\prod_{i\le t}(x-a_i)$, $c=\prod_{i<t}(a_t-a_i)\ne0$ and
$b_i=1/(a_i-a_t)$ for $i<t$; the $b_i$ are distinct and nonzero. The map
$u\mapsto a_t+1/u$ is a bijection from $\Fp^\times$ onto $\Fp\setminus\{a_t\}$,
and $x-a_i=(a_t-a_i)(u-b_i)/u$ for $x=a_t+1/u$ and $i<t$. Hence
\[
 f(a_t+1/u)=\frac{c\,g(u)}{u^t},\qquad g(u)=\prod_{i<t}(u-b_i).
\]
As $t$ is even, $\chi(u^{-t})=1$ for $u\ne0$. Since $f(a_t)=0$ and
$g(0)=\prod_{i<t}(-b_i)=1/c$,
\[
 \sum_{x\in\Fp}\chi(f(x))=\chi(c)\sum_{u\ne0}\chi(g(u))
 =\chi(c)\sum_{u\in\Fp}\chi(g(u))-1 .
\]
The polynomial $g$ is a product of $t-1\ge1$ distinct linear factors, so by
(a) the last sum has modulus at most $(t-2)\sqrt p$.
\end{proof}

Part (b) is the character sum over an even number of distinct roots, with the point at infinity retained: the
substitution moves the root $a_t$ to infinity, and the proof needs only (a).
Its case $t=2$ is the \emph{two-root identity}, which we record with a direct
proof: for distinct $a,b\in\Fp$,
\begin{equation}\label{eq:two-root}
 \sum_{y\in\Fp}\chi\bigl((y-a)(y-b)\bigr)=-1,\qquad\text{hence}\qquad
 \E_{y\in\Fp}\,\chi(y-a)\chi(y-b)=\one_{a=b}-q\quad(a,b\in\Fp).
\end{equation}
Indeed, for $y\ne a$, $\chi((y-a)(y-b))=\chi((y-b)/(y-a))=\chi(1-(b-a)/(y-a))$,
and as $y$ runs over $\Fp\setminus\{a\}$ the argument $1-(b-a)/(y-a)$ runs
over $\Fp\setminus\{1\}$; the term $y=a$ vanishes. So the sum is
$\sum_{s\ne1}\chi(s)=-\chi(1)=-1$, because there are as many squares as
non-squares in $\Fp^\times$. For $a=b$ the average is $(p-1)/p=1-q$.

\begin{theorem}[Gauss sums; see {\cite[Ch.~6, \S\S3--4 and Ch.~8, \S2]{IrelandRosen}}, {\cite[Ch.~1]{BerndtEvansWilliams}}]\label{thm:gauss}
Let $p$ be an odd prime and $\psi$ a nontrivial multiplicative character of $\Fp$.
\begin{enumerate}[label=\textup{(\roman*)}]
\item\label{it:gauss:twist} $g_0(\psi)=0$, and $g_s(\psi)=\bar\psi(s)\,g(\psi)$
  for every $s\in\Fp^\times$.
\item\label{it:gauss:abs} $|g(\psi)|=\sqrt p$.
\item\label{it:gauss:quad} For the quadratic character,
  $\overline{G_\chi}=\chi_p(-1)G_\chi$ and $G_\chi^2=\chi_p(-1)\,p$; moreover
  $G_p(k)=\chi_p(k)G_\chi$ for every $k\in\Fp^\times$, and
  $\sum_{t\in\Fp}\chi_p(t)\ee{kt/p}=\chi_p(k)G_\chi$ for every $k\in\Fp$. We
  write $\eps_G=G_\chi/\sqrt p$; thus $|\eps_G|=1$ and $\eps_G^2=\chi_p(-1)$, so
  $G_\chi$ is real if $p\equiv1\pmod4$ and purely imaginary if $p\equiv3\pmod4$.
\end{enumerate}
\end{theorem}

\begin{proof}
\ref{it:gauss:twist} $g_0(\psi)=\sum_{t\ne0}\psi(t)=0$, because $\psi$ is a
nontrivial character of the finite group $\Fp^\times$: if $\psi(t_0)\ne1$ then
multiplying the sum by $\psi(t_0)$ permutes its terms. For $s\ne0$ substitute
$t=s^{-1}u$: $g_s(\psi)=\sum_u\psi(s^{-1}u)\ee{u/p}=\psi(s^{-1})g(\psi)=\bar\psi(s)g(\psi)$.

\ref{it:gauss:abs} On the one hand, by \ref{it:gauss:twist} and
$|\bar\psi(s)|=1$ for $s\ne0$, $\sum_{s\in\Fp}|g_s(\psi)|^2=(p-1)|g(\psi)|^2$.
On the other hand, expanding,
\[
  \sum_{s\in\Fp}|g_s(\psi)|^2=\sum_{x,y\in\Fp}\psi(x)\overline{\psi(y)}\sum_{s\in\Fp}\ee{s(x-y)/p}=p\sum_{x\in\Fp}|\psi(x)|^2=p(p-1).
\]
Hence $|g(\psi)|^2=p$.

\ref{it:gauss:quad} $\chi_p$ is real, so
$\overline{G_\chi}=\sum_t\chi_p(t)\ee{-t/p}=g_{-1}(\chi_p)=\chi_p(-1)G_\chi$ by
\ref{it:gauss:twist}; with \ref{it:gauss:abs},
$p=G_\chi\overline{G_\chi}=\chi_p(-1)G_\chi^2$. The number of $u\in\Fp$ with
$u^2=t$ is $1+\chi_p(t)$ (this includes $t=0$, since $\chi_p(0)=0$), so for
$k\ne0$
$G_p(k)=\sum_t(1+\chi_p(t))\ee{kt/p}=\sum_t\ee{kt/p}+g_k(\chi_p)=0+\chi_p(k)G_\chi$.
The last formula is \ref{it:gauss:twist} for $k\ne0$ and $g_0(\chi_p)=0$ for
$k=0$. Finally $|\eps_G|=1$ by \ref{it:gauss:abs} and
$\eps_G^2=G_\chi^2/p=\chi_p(-1)$, and $\chi_p(-1)=1$ exactly when
$p\equiv1\pmod4$.
\end{proof}

\begin{remark}[Gauss's sign]\label{rem:gauss-sign}
Gauss evaluated the sign: $G_\chi=\sqrt p$ if $p\equiv1\pmod4$ and
$G_\chi=i\sqrt p$ if $p\equiv3\pmod4$ \cite[Ch.~6, \S4]{IrelandRosen}. It is
not used in this paper: wherever $\eps_G$ occurs, only $|\eps_G|=1$ and
$\eps_G^2=\chi_p(-1)$ are needed, which is \cref{thm:gauss}\ref{it:gauss:quad}.
\end{remark}

\begin{theorem}[Jacobi sums; see {\cite[Ch.~8, \S3]{IrelandRosen}}, {\cite[Ch.~2]{BerndtEvansWilliams}}]\label{thm:jacobi}
Let $p$ be an odd prime and $\psi_1,\psi_2$ multiplicative characters such that
$\psi_1$, $\psi_2$ and $\psi_1\psi_2$ are nontrivial. Then
\[
  \Jac(\psi_1,\psi_2)=\frac{g(\psi_1)\,g(\psi_2)}{g(\psi_1\psi_2)},\qquad\text{hence}\qquad |\Jac(\psi_1,\psi_2)|=\sqrt p .
\]
\end{theorem}

\begin{proof}
Expanding the product and grouping by $s=x+y$,
\[
  g(\psi_1)g(\psi_2)=\sum_{s\in\Fp}\ee{s/p}\sum_{x+y=s}\psi_1(x)\psi_2(y).
\]
For $s=0$ the inner sum is
$\sum_x\psi_1(x)\psi_2(-x)=\psi_2(-1)\sum_{x\ne0}(\psi_1\psi_2)(x)=0$, since
$\psi_1\psi_2$ is nontrivial (\cref{thm:gauss}\ref{it:gauss:twist}, first part
of its proof). For $s\ne0$ substitute $x=st$, $y=s(1-t)$; this is a bijection
of the solutions of $x+y=s$ with $t\in\Fp$, and the inner sum becomes
$(\psi_1\psi_2)(s)\sum_t\psi_1(t)\psi_2(1-t)=(\psi_1\psi_2)(s)\Jac(\psi_1,\psi_2)$.
Hence
\[
 g(\psi_1)g(\psi_2)=\Jac(\psi_1,\psi_2)\sum_{s\ne0}(\psi_1\psi_2)(s)\ee{s/p}
 =\Jac(\psi_1,\psi_2)\,g(\psi_1\psi_2).
\]
By \cref{thm:gauss}\ref{it:gauss:abs} all three Gauss sums have modulus
$\sqrt p\ne0$, which gives both claims.
\end{proof}

\subsection{The square shift and the Paley matrix}\label{subsec:notation:square-shift}

The square shift adds a uniformly chosen nonzero square. At the primes $p\equiv3\pmod4$ dividing
the small modulus $\Mzero$ it is the one-coordinate kernel of the external index
(\cref{def:hyb:coupled}), so its centred norm is the inclusion probability of
that coordinate; its tensor products over a set of primes give the
strict-square pair law, which the
cancellation along an accepted query expands in characters; and the Paley
matrix is the signed part of both.

\begin{lemma}[The square shift and the Paley matrix]\label{lem:square-shift}
Let $p$ be an odd prime, $q=1/p$ and $r=p^{-1/2}$. Let $T_p$ be the
\emph{square shift} on $L^2(\Fp)$,
\[
 (T_pf)(x)=\E_{u\in\Qsq}f(x+u),
\]
let $C$ be the $p\times p$ \emph{Paley matrix} $C(x,y)=\chi(y-x)$, and let
$P=T_C$ be the \emph{Paley operator}, $(Pf)(x)=\E_{y\in\Fp}\chi(y-x)f(y)$.
\begin{enumerate}[label=\textup{(\alph*)}]
\item $T_p$ is doubly stochastic: its matrix $\frac2{p-1}\Ap(y-x)$ has
 $(p-1)/2$ nonzero entries in each row and in each column. It is diagonal in
 the additive characters: $T_pe_0=e_0$, and $T_pe_k=\lambda_ke_k$ for $k\ne0$
 with
 \[
  \lambda_k=\frac{G_p(k)-1}{p-1}=\frac{\chi(k)G_\chi-1}{p-1}.
 \]
 Hence $\|T_p|_{L^2_0}\|\le(\sqrt p+1)/(p-1)=1/(\sqrt p-1)$, with equality if
 $p\equiv1\pmod4$. Likewise the matrices $(\Ap(\pm(y-x))-\frac12)_{x,y}$ have
 the eigenvalue $-\frac12$ at $e_0$ and $(G_p(\pm k)-1)/2$ at $e_k$, $k\ne0$,
 and norm at most $(\sqrt p+1)/2$.
\item If $p\equiv3\pmod4$, then $|\lambda_k|=\sqrt{p+1}/(p-1)$ for every
 $k\ne0$, so $T_p$ restricted to $L^2_0(\Fp)$ is $\sqrt{p+1}/(p-1)$ times a
 unitary operator, and
 \[
  \|T_p|_{L^2_0}\|=\frac{\sqrt{p+1}}{p-1}.
 \]
 The norm in this display is best possible among such transitions: for every odd $p$, every doubly stochastic matrix $T$ on
 $\Fp$ with $T(x,y)=0$ unless $y-x\in\Qsq$ satisfies
 $\|T|_{L^2_0}\|\ge\sqrt{p+1}/(p-1)$. The number $\sqrt{p+1}/(p-1)$ is $1$ at
 $p=3$, less than $1$ for $p\ge5$, and less than $1/(\sqrt p-1)$ for every $p$.
\item $C\one=C^*\one=0$, $C^*=\chi(-1)C$ and $CC^*=C^*C=pI-\Jall$, so
 $\|C\|=\sqrt p$. Hence $P\one=P^*\one=0$, $P^*=\chi(-1)P$,
 $P^*P=PP^*=q(I-\Jproj)$ and $\|P\|=r$. Moreover, as matrices and as
 operators on $L^2(\Fp)$,
 \[
  T_p=\frac{\Jall-I+C}{p-1},\qquad T_p=\frac{p\Jproj-I+pP}{p-1}.
 \]
 In particular $C^*=-C$ and $P^*=-P$ if $p\equiv3\pmod4$.
\item Let $\mathcal S$ be a finite set of odd primes,
 $X_{\mathcal S}=\prod_{p\in\mathcal S}\Fp$ with the uniform probability, and
 $\pi_{\mathcal S}$ the uniform probability on
 $\prod_{p\in\mathcal S}\Om\subseteq X_{\mathcal S}\times X_{\mathcal S}$, the
 strict-square pair law. The characters of $X_{\mathcal S}$ are
 $z\mapsto\ee{z\xi}=\prod_p\ee{a_pz_p/p}$ with $\xi=\sum_pa_p/p\in\Q/\Z$, and
 $\widehat g(\xi)=\E_zg(z)\ee{-z\xi}$. For functions $g_1,g_2$ on
 $X_{\mathcal S}$,
 \[
  \int g_1(x)g_2(y)\,d\pi_{\mathcal S}(x,y)=\sum_{\xi}\widehat{g_1}(-\xi)\,\widehat{g_2}(\xi)\,g(\xi),
 \]
 the sum running over the characters of $X_{\mathcal S}$, where
 \[
  g(\xi)=\prod_{p\in\mathcal S,\ a_p\ne0}\frac{G_p(a_p)-1}{p-1}
  =\frac1{\varphi(q(\xi))}\sum_{m\in(\Z/q(\xi)\Z)^\times}\ee{\xi m^2}
 \]
 is the unit quadratic average of the actual-square kernel
 (\cref{def:kernel-notation}). Moreover
 \[
  |g(\xi)|\le\prod_{p\mid q(\xi)}\frac1{\sqrt p-1},\qquad\text{and}\qquad
  |g(\xi)|\le q(\xi)^{-1/3}\ \text{if every prime factor of $q(\xi)$ is at least $64$.}
 \]
\end{enumerate}
\end{lemma}

\begin{proof}
(a) Since $|\Qsq|=(p-1)/2$, the matrix of $T_p$ is
$T_p(x,y)=\frac2{p-1}\Ap(y-x)$. For fixed $x$ there are $(p-1)/2$ values of $y$
with $y-x\in\Qsq$, and for fixed $y$ there are as many values of $x$; so $T_p$
is doubly stochastic. It is a convolution, and
$T_pe_k(x)=e_k(x)\,\E_{u\in\Qsq}\ee{ku/p}$. Every nonzero square has exactly two
square roots, so
\[
 \E_{u\in\Qsq}\ee{ku/p}=\frac1{p-1}\sum_{m\in\Fp^\times}\ee{km^2/p}=\frac{G_p(k)-1}{p-1}.
\]
This is $1$ at $k=0$, and $(\chi(k)G_\chi-1)/(p-1)$ at $k\ne0$ by
\cref{thm:gauss}\ref{it:gauss:quad}. The $e_k$ with $k\ne0$ form an orthonormal
basis of $L^2_0(\Fp)$, so $\|T_p|_{L^2_0}\|=\max_{k\ne0}|\lambda_k|\le
(|G_\chi|+1)/(p-1)=(\sqrt p+1)/(p-1)$. If $p\equiv1\pmod4$, then
$G_\chi=\pm\sqrt p$ is real, $\chi$ takes both signs on $\Fp^\times$, and some
$k\ne0$ has $\chi(k)G_\chi=-\sqrt p$, so that $|\lambda_k|=(\sqrt p+1)/(p-1)$.
For the second statement,
$\sum_y(\Ap(y-x)-\frac12)e_k(y)=e_k(x)\bigl(\sum_{u\in\Qsq}\ee{ku/p}-\frac p2\one_{k=0}\bigr)$,
which is $-\frac12e_0$ at $k=0$ and $\frac12(G_p(k)-1)e_k$ at $k\ne0$; for
$\Ap(x-y)$ replace $k$ by $-k$. Both matrices are diagonal in the orthonormal
basis $(p^{-1/2}e_k)_k$ of $\C^{\Fp}$, so their norms are the largest moduli of
their eigenvalues, at most $(\sqrt p+1)/2$, since $|G_p(\pm k)|=\sqrt p$ for
$k\ne0$.

(b) If $p\equiv3\pmod4$, then $G_\chi^2=-p$ (\cref{thm:gauss}\ref{it:gauss:quad}),
so $\chi(k)G_\chi=\pm i\sqrt p$ for $k\ne0$ and $|\chi(k)G_\chi-1|^2=p+1$.
Hence $|\lambda_k|=\sqrt{p+1}/(p-1)$ for every $k\ne0$, and $T_p$ is diagonal on
$L^2_0(\Fp)$ in an orthonormal basis with eigenvalues of this common modulus.
For the optimality, let $T$ be doubly stochastic with $T(x,y)=0$ unless
$y-x\in\Qsq$. Since $T\one=\one$ and $T^*\one=\one$, $T$ maps $L^2_0$ into
itself, and the constants carry the singular value $1$. Hence the squared
Hilbert--Schmidt norm satisfies
$\|T\|_{\rm HS}^2=1+\|T|_{L^2_0}\|_{\rm HS}^2\le1+(p-1)\|T|_{L^2_0}\|^2$. Each
row of $T$ has at most $(p-1)/2$ nonzero entries, summing to $1$, so by
Cauchy--Schwarz its squared entries sum to at least $2/(p-1)$, and
$\|T\|_{\rm HS}^2\ge2p/(p-1)$. Combining,
$\|T|_{L^2_0}\|^2\ge(2p/(p-1)-1)/(p-1)=(p+1)/(p-1)^2$. Finally
$\sqrt{p+1}<\sqrt p+1$ for every $p$, and $\sqrt{p+1}<p-1$ exactly when $p>3$.

(c) The row sums $\sum_y\chi(y-x)$ and the column sums vanish, so
$C\one=C^*\one=0$. The adjoint has entries
$C^*(x,y)=\overline{C(y,x)}=\chi(x-y)=\chi(-1)C(x,y)$. The $(x,z)$ entry of $CC^*$
is $\sum_y\chi(y-x)\chi(y-z)$, which is $p-1$ for $x=z$ and $-1$ for $x\ne z$ by
\eqref{eq:two-root}; so $CC^*=pI-\Jall$, and $C^*C=\chi(-1)C^2=CC^*$. On the
constants $CC^*$ vanishes and on their orthogonal complement it is $pI$, so
$\|C\|^2=p$. As a matrix acting on vectors, $P$ is $C/p$, and the uniform
probability multiplies every inner product by the same factor $1/p$, so the
adjoint of $P$ is $C^*/p$. Hence $P^*=\chi(-1)P$,
$P^*P=C^*C/p^2=(pI-\Jall)/p^2=q(I-\Jproj)$ (as $\Jall/p=\Jproj$ on vectors),
$PP^*=P^*P$ and $\|P\|=\|C\|/p=r$. Finally
$2\Ap(y-x)=1+\chi(y-x)-\one_{y=x}$ for all $x,y$: both sides are $2$ if
$y-x\in\Qsq$ and $0$ otherwise. So
$T_p=\frac2{p-1}(\Ap(y-x))_{x,y}=(\Jall-I+C)/(p-1)$; applied to $f$ this is
$(p\E f-f+pPf)/(p-1)$.

(d) The map $(x,u)\mapsto(x,x+u)$ is a bijection of $\Fp\times\Qsq$ onto $\Om$.
So under $\pi_{\mathcal S}$ the pair $(x,y)$ is $(x,x+u)$, with $x$ uniform on
$X_{\mathcal S}$ and $u=(u_p)_p$ independent of $x$, its coordinates independent
and uniform on the sets $\Qsq$. Hence
$\int g_1(x)g_2(y)\,d\pi_{\mathcal S}=\E_xg_1(x)(T_{\mathcal S}g_2)(x)$ with
$T_{\mathcal S}=\bigotimes_{p\in\mathcal S}T_p$. Expand
$g_2(y)=\sum_\xi\widehat{g_2}(\xi)\ee{y\xi}$. The character $\ee{\cdot\,\xi}$ is
the tensor product of the characters $e_{a_p}$ of the factors, so by (a)
$T_{\mathcal S}\ee{\cdot\,\xi}=g(\xi)\ee{\cdot\,\xi}$ with
$g(\xi)=\prod_{a_p\ne0}(G_p(a_p)-1)/(p-1)$; and
$\E_xg_1(x)\ee{x\xi}=\widehat{g_1}(-\xi)$. This is the identity. By the
computation in (a), $(G_p(a)-1)/(p-1)=\frac1{p-1}\sum_{m\in\Fp^\times}\ee{am^2/p}$.
By the Chinese remainder theorem, $m\bmod q(\xi)$ runs over the units exactly
when $(m\bmod p)_{p\mid q(\xi)}$ runs over $\prod_p\Fp^\times$, and
$\ee{\xi m^2}=\prod_p\ee{a_pm^2/p}$; this gives the second expression for
$g(\xi)$. By (a) each factor has modulus at most $1/(\sqrt p-1)$, and for
$p\ge64$, $\sqrt p-1=p^{1/3}(p^{1/6}-p^{-1/3})\ge p^{1/3}(2-\tfrac14)\ge p^{1/3}$.
Multiply over $p\mid q(\xi)$.
\end{proof}

The optimality in (b) concerns this specified class of transitions, the
doubly stochastic kernels supported on square shifts; it places no
restriction on a different construction of the indices. The index law uses
(b), the Paley channels of the local laws use (c) (\cref{lem:paley-triple}),
the pinned pair bounds (bounds with prescribed labels fixed) use
the second statement of (a) (\cref{lem:mid:pinpair}), and the
cancellation uses (d).

\subsection{Primes in the classes modulo four}\label{subsec:notation:primes}

For primes in a residue class modulo four we use the following explicit
Chebyshev estimate. Write $\theta(x;4,r)=\sum_{p\le x,\ p\equiv r\ (4)}\log p$.

\begin{theorem}[Bennett--Martin--O'Bryant--Rechnitzer {\cite[Theorem~1.2 and Table~1]{BMOR18}}]\label{thm:bmor}
For $r\in\{1,3\}$ and every real $x\ge4\,800\,162\,889$,
\[
 \Bigl|\theta(x;4,r)-\frac x2\Bigr|<c_\theta(4)\,\frac{x}{\log x},
 \qquad c_\theta(4)=0.0004822=\frac{2411}{5000000}.
\]
Consequently, if $X\ge4\,800\,162\,889$, then
$|\theta(t;4,r)-t/2|\le\mathsf e_Xt$ for every $t\ge X$, where
$\mathsf e_X=2411/(5000000\log X)$.
\end{theorem}

The consequence holds because $t/\log t\le t/\log X$ for $t\ge X$. The constants
are those of the published version (\cref{subsec:external:sources}). In the
construction this theorem supplies the prime counts for the cells of the middle range above
$4\,800\,162\,890$, from which their prices and the analytic lower start sums of the mean credit are computed
(\cref{eq:record:cell-charge,lem:record:lower-start-sums}); the same cell charge
prices the cells of the upper bands (\cref{prop:record:upper}). Below
$4\,800\,162\,890$, one more than the threshold of the theorem, no Chebyshev estimate is used (\cref{rem:record:middle-scope}).

\subsection{Two correlation inequalities (Peyre and Kassabov)}\label{subsec:notation:correlation}

Maximal correlation and the Friedrichs angle are defined in
\cref{def:prelim:maxcorr,def:prelim:angle}. When the intersection of the two
function spaces consists of the constants, the maximal correlation is the
cosine of the Friedrichs angle between them; this is why an inequality for
angles (Kassabov's) and one for correlations (Peyre's) combine in the pinned
influence bounds of the local laws.

\begin{theorem}[Peyre's tensorization {\cite[Theorem~3.3.1]{Peyre}}]\label{thm:peyre-finite}
Let $A,B$ be disjoint nonempty finite sets, and let $(X_v)_{v\in A\cup B}$ be
random variables on a common probability space, each taking finitely many
values. Let $\eps_{ij}\ge0$ for $i\in A$, $j\in B$, and
$\mathcal E=(\eps_{ij})$. Suppose that for all $i\in A$, $j\in B$, every
$S\subseteq(A\cup B)\setminus\{i,j\}$ and every value $s$ of
$X_S=(X_v)_{v\in S}$ with $\Pr(X_S=s)>0$,
\begin{equation}\label{eq:peyre:hyp}
  \rho(X_i;X_j\mid X_S=s)\le\eps_{ij}.
\end{equation}
Then
\begin{equation}\label{eq:peyre:concl}
  \rho\big((X_a)_{a\in A};(X_b)_{b\in B}\big)\le\min\{1,\opnorm{\mathcal E}\}.
\end{equation}
More precisely, \eqref{eq:peyre:concl} holds as soon as there are enumerations
$A=\{a_1,\dots,a_N\}$ and $B=\{b_1,\dots,b_M\}$ such that \eqref{eq:peyre:hyp}
holds for $i=a_k$, $j=b_l$ and the single set
$S=\{a_1,\dots,a_{k-1}\}\cup\{b_1,\dots,b_{l-1}\}$, for all $1\le k\le N$ and
$1\le l\le M$.
\end{theorem}

Here $\opnorm{\mathcal E}$ is the norm of $\mathcal E$ as an operator
$\ell^2(B)\to\ell^2(A)$ (Euclidean norms, \cref{def:prelim:operators}), its
largest singular value. The second statement has a weaker hypothesis, because
the sets $S$ it uses are among those allowed in the first: they avoid $a_k$
and $b_l$. \Cref{app:peyre} gives the complete proof, which is Peyre's argument in
finite form; the proof establishes the second statement directly. The paper
applies only the first.

\begin{corollary}[Independent pairs]\label{cor:peyre-independent}
Let $(X_k,Y_k)$, $1\le k\le n$, be pairs of finite-valued random variables,
the $n$ pairs independent of each other. Then
\[
 \rho\big((X_k)_k;(Y_k)_k\big)=\max_k\rho(X_k;Y_k).
\]
The same holds conditionally: if the triples $(X_k,Y_k,Z_k)$ are independent
across $k$ and $z=(z_k)_k$ is a value of $(Z_k)_k$ of positive probability,
then $\rho((X_k)_k;(Y_k)_k\mid(Z_k)_k=z)=\max_k\rho(X_k;Y_k\mid Z_k=z_k)$.
\end{corollary}

\begin{proof}
Apply \cref{thm:peyre-finite} with $X_{a_k}=X_k$ and $X_{b_k}=Y_k$. Let $S$
avoid $a_k$ and $b_l$, and let $s$ be a value of $X_S$ of positive probability.
The event $\{X_S=s\}$ is an intersection of events, one in each pair, so under
$\Pr(\cdot\mid X_S=s)$ the pairs are still independent. If $k\ne l$, then
$X_k$ and $Y_l$ lie in different pairs, are conditionally independent, and
$\rho(X_k;Y_l\mid X_S=s)=0$. If $k=l$, the event depends only on the other
pairs, the conditional law of $(X_k,Y_k)$ is its law, and
$\rho(X_k;Y_k\mid X_S=s)=\rho(X_k;Y_k)$. So the hypothesis holds with the
diagonal matrix $\mathcal E=\diag(\rho(X_k;Y_k))_k$, whose norm is
$\max_k\rho(X_k;Y_k)$. The reverse inequality holds because functions of $X_k$
alone and of $Y_k$ alone are admissible in the definition of
$\rho((X_k)_k;(Y_k)_k)$. For the conditional statement apply this under
$\Pr(\cdot\mid(Z_k)_k=z)$: by the same product computation the pairs remain
independent, with the conditional laws of $(X_k,Y_k)$ given $Z_k=z_k$.
\end{proof}

In operator form: the operator of a product of couplings is the tensor product
of the operators of the factors, and its centred norm is the largest centred
norm of a factor (\cref{lem:sector}(c)).

\begin{corollary}[Conditional decorrelation]\label{cor:cond-decorrelation}
Let $X,Y,Z$ be finite-valued random variables on a common probability space
and $\eps\ge0$, and assume $\rho(X;Y\mid Z=z)\le\eps$ for every value $z$ with
$\Pr(Z=z)>0$. If a real function $f\in L^2(X,Z)$ satisfies $\E[f\mid Z]=0$ and
$g\in L^2(Y,Z)$ is real, then
\[
 |\E[fg]|\le\eps\,\norm f\,\norm{g-\E[g\mid Z]}\le\eps\norm f\norm g .
\]
\end{corollary}

\begin{proof}
This is \cref{lem:peyre:basic}(b), proved in \cref{app:peyre} on a finite space
of full support (\cref{sec:peyre:conv}); every quantity here depends only on
the joint law of $(X,Y,Z)$, so this is no restriction. The proof is fibrewise:
on each set $\{Z=z\}$ of positive probability the hypothesis bounds
$|\E[fg\mid Z=z]|$ by $\eps$ times the conditional norms of $f$ and of
$g-\E[g\mid Z]$, and Cauchy--Schwarz for the sum over $z$ gives the claim.
\end{proof}

The two corollaries are used together when a law on a whole block
is built from feature functions (the block step of
\cref{lem:oct24:feature-block}): given a function of all the
labels of a block, the conditional correlation tensorizes over
the primes of the block
(\cref{cor:peyre-independent}), and a function with conditional mean zero is
then decorrelated from every function of the other side
(\cref{cor:cond-decorrelation}).

\begin{theorem}[Kassabov's angle inequality {\cite[Lemma~4.2(b)]{Kassabov}}]\label{thm:kassabov}
Let $V_1,V_2,V_3$ be subspaces of a finite-dimensional Hilbert space and put
$\eps_{ij}=\cos\angle(V_i,V_j)$. If $\eps_{13}<1$ and $\eps_{23}<1$, then
\[
  \cos\angle(V_1\cap V_3,V_2\cap V_3)\le\frac{\eps_{12}+\eps_{13}\eps_{23}}{\sqrt{(1-\eps_{13}^2)(1-\eps_{23}^2)}} .
\]
\end{theorem}

It enters through the elimination of pinned vertices in \cref{lem:mid:allpin},
and through it in the pinned influence theorem \cref{prop:middle-negative}.

\subsection{Monotone majorants: a bound at the smallest prime implies the bound at every larger prime}\label{subsec:notation:majorants}

Every local constant of Part~V is an explicit expression in a prime $p$. It is
computed once, at the start $X$ of a range of primes (a piece of the middle
range, an upper band, the tail threshold), and then used at every prime
$p\ge X$ and on every block of primes of the range. This is valid because each
such constant, written as a function of the real variables $r=p^{-1/2}$,
$q=1/p$, $\lambda=1/\log p$ and a block count (the number of primes in the block), is nondecreasing in each of
them. To write a constant in these variables replace $p$ by $r^{-2}$: for
instance $p/(p-1)=1/(1-r^2)$, $1/(p-1)=r^2/(1-r^2)$, $k/(p-k)=kr^2/(1-kr^2)$ and
$k/p=kr^2$ for a fixed number $k$. ``Nondecreasing'' always refers to these variables. The definition
and the lemma below record the calculus once; every section that evaluates a
constant at an endpoint cites the lemma and displays the constant in a form
built from its operations.

\begin{definition}[Admissible and small functions]\label{def:majorant}
A \emph{box} is a product $\Delta=I_1\times\dots\times I_k$ of intervals
$I_j\subseteq[0,\infty)$, each containing its supremum $x^+_j<\infty$; the
point $x^+=(x^+_1,\dots,x^+_k)$ is its \emph{endpoint}, and $\Delta$ is ordered
coordinatewise. A function $\varphi\colon\Delta\to[0,\infty)$ is
\emph{admissible} if $\varphi(x)\le\varphi(x')$ whenever $x\le x'$. If $r$ is a
coordinate whose interval lies in $(0,\infty)$, then $\varphi$ is \emph{small} if
$\varphi/r$ is admissible. A map $F\colon\mathcal D\to[0,\infty)$ on a set
$\mathcal D\subseteq[0,\infty)^m$ is \emph{monotone} if $F(y)\le F(y')$ whenever
$y\le y'$ in $\mathcal D$.
\end{definition}

\begin{lemma}[Admissible majorants]\label{lem:majorant}
Let $\Delta$ be a box with endpoint $x^+$.
\begin{enumerate}[label=\textup{(\roman*)}]
\item Nonnegative constants, the coordinates, and every power series in the
 coordinates with nonnegative coefficients that converges on $\Delta$ are
 admissible; such a series in which every monomial contains $r$ is small. In
 particular every polynomial in $r$ with nonnegative coefficients, such as a
 count error of \cref{lem:extension}, and $p/(p-1)=1/(1-r^2)$ are admissible.
\item If $F\colon\mathcal D\to[0,\infty)$ is monotone and
 $\varphi_1,\dots,\varphi_m$ are admissible with
 $(\varphi_1(x),\dots,\varphi_m(x))\in\mathcal D$ for every $x\in\Delta$, then
 $F(\varphi_1,\dots,\varphi_m)$ is admissible. Hence sums, products, maxima,
 minima and square roots of admissible functions, and powers $\varphi^\psi$
 with $\varphi\ge1$, are admissible; so are $e^\varphi$, $e^\varphi-1$,
 $(e^\varphi-1)/\varphi$ (with value $1$ where $\varphi=0$) and
 $\log(1+\varphi)$, and, where $\varphi<1$ on $\Delta$, also $1/(1-\varphi)$,
 $(1+\varphi)/(1-\varphi)$ and $-\log(1-\varphi)$; and so are the operator
 norm of a matrix whose entries are admissible and the largest eigenvalue of
 a real symmetric matrix whose entries are admissible. Finally,
 $(G,D,r)\mapsto G+D^2r/(1+Gr)$ is monotone on
 $\{G,D,r\ge0,\ Dr\le1\}$, so $\varphi_G+\varphi_D^2\varphi_r/(1+\varphi_G\varphi_r)$
 is admissible for admissible $\varphi_G,\varphi_D,\varphi_r$ with
 $\varphi_D\varphi_r\le1$ on $\Delta$.
\item Every small function is admissible. If $\varphi,\psi$ are small and
 $h$ is admissible, then $\varphi+\psi$, $\max(\varphi,\psi)$,
 $\min(\varphi,\psi)$, $h\varphi$ and $\sqrt{\varphi\psi}$ are small. If
 $F(y)=yG(y)$ with $G\ge0$ monotone, then $F\circ\varphi$ is small for small
 $\varphi$ (with values in the domain of $G$); in particular $e^\varphi-1$,
 and, where $\varphi<1$, $-\log(1-\varphi)$, $\varphi/(1-\varphi)$ and
 $(1+\varphi)/(1-\varphi)-1$ are small. If $\varphi_1,\dots,\varphi_n$ are
 small, then $\prod_i(1+\varphi_i)-1$ is small. The function $\log(1+\varphi)$
 need not be small: $\log(1+y)/y$ decreases.
\item One fixed candidate, chosen independently of $x$ (for instance by
 its value at the endpoint), among finitely many admissible candidates is admissible. If $\varphi$ is admissible and $\psi\le\varphi$ on
 $\Delta$, then $\max(\varphi,\psi)=\varphi$ is admissible, whatever $\psi$ is.
\end{enumerate}
Consequently the following hold.
\begin{enumerate}[label=\textup{(\alph*)}]
\item \emph{Endpoint.} If $\varphi$ is admissible, then $\varphi(x)\le\varphi(x^+)$
 for every $x\in\Delta$. In particular, let $P>1$ and $L\ge0$ be real, and let
 $\Delta=(0,P^{-1/2}]\times(0,1/P]\times(0,1/\log P]\times[0,L/\log P]$ in the
 variables $r$, $q$, $\lambda$ and the block count (or a product of some of
 these factors). Then every inequality of the form $\varphi(x^+)<c$ (or $\le c$) checked at the
 endpoint holds at every prime $p\ge P$ and every block count at most
 $L/\log P$, and every admissible coefficient at such a prime is at most its
 value at the endpoint.
\item \emph{Piece start.} Let $X>1$, $q_X=1/X$, $r_X\ge X^{-1/2}$ and
 $n_X\ge0$ (the certificates use $r_X=1/\lfloor\sqrt X\rfloor$), and let
 $\varphi(q,r,l)$ be admissible on the box $(0,q_X]\times(0,r_X]\times[0,n_X]$.
 If a quantity at the prime $p$ with block count $l\le n_X$ is
 $r^d\varphi(q,r,l)$, where $d\ge0$, $r=p^{-1/2}$ and $q=1/p$, then for every
 prime $p\ge X$ it is at most $p^{-d/2}\varphi(q_X,r_X,n_X)$.
\item \emph{Block inflation.} Let $c\ge0$ and $L>0$. The function
 $X\mapsto\exp(cL/(X\log X))$ is nonincreasing on $(1,\infty)$ (decreasing if
 $c>0$). If $\Blk$ is a set of
 primes, all at least $X>1$, with $q_\Blk\le e^L$, then $|\Blk|\le L/\log X$,
 $\sum_{p\in\Blk}1/p\le L/(X\log X)$ and
 $\prod_{p\in\Blk}(1+c/p)\le\exp(cL/(X\log X))$.
\item \emph{Matrix tests.} Let $M_0$ be a fixed Hermitian matrix and
 $E_1,\dots,E_m\succeq0$ fixed positive semidefinite matrices of the same
 size, and let $c_1,\dots,c_m$ be real functions on $\Delta$
 with $c_i(x)\ge c_i(x^+)$ for every $x$ (for instance $c_i=-\varphi_i$, or
 $c_i=g/\varphi_i$ with a constant $g\ge0$ independent of $x$, where $\varphi_i>0$ is admissible). If
 $M_0+\sum_ic_i(x^+)E_i\succeq0$, then $M_0+\sum_ic_i(x)E_i\succeq0$ for every
 $x\in\Delta$.
\item \emph{Several thresholds.} Let $\mathcal A(p,g)$ be a statement about a
 prime $p$ and a number $g$ that remains true when $g$ is increased (for
 instance $C_p\succeq-gp^{-1/2}I$, $C_p\preceq(1+g/p)I$ or $c(p)\le g$). If,
 for each $s$ in a finite set, $\mathcal A(p,g_s)$ holds for every prime
 $p\ge y_s$, then $\mathcal A(p,\min\{g_s:y_s\le y\})$ holds for every prime
 $p\ge y$ and every real $y\ge\min_sy_s$.
\end{enumerate}
\end{lemma}

\begin{proof}
(i) For $0\le x\le x'$ every monomial satisfies $x^\alpha\le x'^\alpha$, so a
series with nonnegative coefficients is admissible. If every monomial contains
$r$, the series is $r$ times a series with nonnegative coefficients, which
converges on $\Delta$ because $r>0$ there; so it is small. The count errors are
polynomials in $r$ with nonnegative coefficients, and
$1/(1-r^2)=\sum_{k\ge0}r^{2k}$ for $r<1$.

(ii) If $x\le x'$, then $\varphi_i(x)\le\varphi_i(x')$ for every $i$, and the
monotonicity of $F$ gives the claim. The listed maps are monotone on their
domains: $(y,z)\mapsto y+z$, $yz$, $\max(y,z)$, $\min(y,z)$ on $[0,\infty)^2$;
$y\mapsto\sqrt y$; $(y,z)\mapsto y^z$ on $[1,\infty)\times[0,\infty)$;
$y\mapsto e^y$, $e^y-1$, $(e^y-1)/y=\sum_{k\ge0}y^k/(k+1)!$ and $\log(1+y)$ on
$[0,\infty)$; and $y\mapsto1/(1-y)=\sum_ky^k$,
$(1+y)/(1-y)=1+2\sum_{k\ge1}y^k$ and $-\log(1-y)=\sum_{k\ge1}y^k/k$ on $[0,1)$.
For real matrices with $0\le M\le M'$ entrywise and unit vectors $u,v$,
$|v^*Mu|\le|v|^{\mathsf T}M|u|\le|v|^{\mathsf T}M'|u|\le\|M'\|$, where $|u|$ is
the vector of moduli; so $\|M\|\le\|M'\|$. For real symmetric $M\ge0$,
$u^{\mathsf T}Mu\le|u|^{\mathsf T}M|u|$, so the largest eigenvalue is the maximum
of $u^{\mathsf T}Mu$ over nonnegative unit vectors $u$, which is monotone in the
entries in the same way. Finally, on
$\mathcal D=\{G,D,r\ge0,\ Dr\le1\}$ the map $F(G,D,r)=G+D^2r/(1+Gr)$ has the
partial derivatives $1-D^2r^2/(1+Gr)^2\ge0$ (as $Dr\le1\le1+Gr$),
$2Dr/(1+Gr)\ge0$ and $D^2/(1+Gr)^2\ge0$. If $y\le y'$ in $\mathcal D$, the
path from $y$ to $y'$ that raises first $G$, then $D$, then $r$ stays in
$\mathcal D$, because the product $Dr$ only increases along it, up to its
value at $y'$; so $F(y)\le F(y')$.

(iii) A small $\varphi=r\varphi'$ is a product of admissible functions. Write
also $\psi=r\psi'$. Then $\varphi+\psi=r(\varphi'+\psi')$,
$\max(\varphi,\psi)=r\max(\varphi',\psi')$,
$\min(\varphi,\psi)=r\min(\varphi',\psi')$, $h\varphi=rh\varphi'$ and
$\sqrt{\varphi\psi}=r\sqrt{\varphi'\psi'}$, and (ii) applies to the cofactors.
If $F(y)=yG(y)$, then $F(\varphi)/r=G(\varphi)\varphi'$, admissible by (ii);
this covers $e^y-1$ with $G(y)=(e^y-1)/y$, $-\log(1-y)$ with
$G(y)=\sum_{k\ge1}y^{k-1}/k$, $y/(1-y)$ with $G(y)=1/(1-y)$, and
$(1+y)/(1-y)-1=2y/(1-y)$. Finally
$\prod_i(1+\varphi_i)-1=\sum_{\emptyset\ne I}\prod_{i\in I}\varphi_i$, and each
term is a small function times admissible ones.

(iv) A choice made once is a fixed admissible function, and
$\max(\varphi,\psi)=\varphi$ when $\psi\le\varphi$.

(a) Every $x\in\Delta$ satisfies $x\le x^+$. For a prime $p\ge P>1$,
$p^{-1/2}\le P^{-1/2}$, $1/p\le1/P$ and $1/\log p\le1/\log P$.

(b) For a prime $p\ge X$, $1/p\le q_X$, $p^{-1/2}\le X^{-1/2}\le r_X$ and
$l\le n_X$, so $\varphi(1/p,p^{-1/2},l)\le\varphi(q_X,r_X,n_X)$ by (a); multiply
by $p^{-d/2}$.

(c) The function $X\log X$ is positive on $(1,\infty)$, with derivative
$\log X+1>0$. Every prime of $\Blk$ is at least $X$, so
$|\Blk|\log X\le\log q_\Blk\le L$; then $\sum_{p\in\Blk}1/p\le|\Blk|/X\le
L/(X\log X)$, and $\prod_{p\in\Blk}(1+c/p)\le\exp(c\sum_{p\in\Blk}1/p)$.

(d) The difference of the two matrices is
$\sum_i(c_i(x)-c_i(x^+))E_i\succeq0$.

(e) Every $s$ with $y_s\le y$ gives $\mathcal A(p,g_s)$ at every prime
$p\ge y\ge y_s$, hence $\mathcal A(p,g)$ for the smallest such $g_s$.
\end{proof}

\begin{remark}[How the lemma is applied]\label{rem:majorant-use}
(1) \emph{Cancel the powers of $r$ first.} A remainder of the form
$q\cdot c(q,r)$ is used as $r\cdot(rc(q,r))$, with the admissible coefficient
$rc(q,r)$, before the piece start is used. At the piece start the coefficient
is evaluated at $q_X=1/X$ and $r_X=1/\lfloor\sqrt X\rfloor\ge X^{-1/2}$, and the bound used for a
block with least prime $b\ge X$ is $b^{-1/2}$ times this
coefficient (part (b) with $d=1$). A remainder computed at $q_X$ and then
divided by the larger $r_X$ would not in general bound the coefficient, since
$q_X/r_X<X^{-1/2}$ when $r_X>X^{-1/2}$; this is never done.

(2) \emph{Logarithmic counts.} Inserting the vertices of a labelling in a
fixed order, with the count errors of \cref{lem:extension} at each step,
gives logarithmic count bounds. Where such a bound must be small in the sense
of \cref{def:majorant} it is used in the form $\sum e$, an
upper bound for $\sum\log(1+e)$, or $\sum-\log(1-e)$; both are small when the
count errors $e$ are, whereas $\sum\log(1+e)$ is admissible but in general not
small.

(3) \emph{Choices and gates.} Orders, partitions, Schur-test parameters and the
branch of a maximum are fixed at the endpoint before (a) is applied. A minimum
of bounds is taken only among bounds whose own hypotheses, which are gates,
hold at the endpoint, and hence everywhere by (a). No monotonicity of an
eigenvector or of an optimized parameter is used. When one branch of a maximum
is not monotone, (iv) applies once it is shown to be dominated by an admissible
branch; this is how the two branches of the maximum in the spectral criterion
are treated (\cref{lem:spectral-monotone,lem:record:monotone-endpoint}).

(4) \emph{Classes.} A matrix test whose matrices depend on $p$ only through
$p\bmod4$ is checked for both classes, and (d) is applied within each class.
The threshold hypotheses of spectral data are gates built from the
operations of (i) and (ii), and persist by (a). A spectral cap valid for every
prime $p\ge P$, times the block inflation of (c) evaluated at the start
$X\ge P$ of a piece or band, bounds the corresponding factor on every block of
that piece or band (\cref{cor:spectral-criterion-blocks}).
\end{remark}

\paragraph{Consequences and what to read next.}
External inputs: Weil's theorem (\cref{thm:weil}(a)), the explicit Chebyshev
bound of Bennett, Martin, O'Bryant and Rechnitzer (\cref{thm:bmor}) and
Kassabov's inequality (\cref{thm:kassabov}); \cref{subsec:external:sources}
gives their sources. Peyre's theorem (\cref{thm:peyre-finite}) and the
conditional decorrelation inequality (\cref{cor:cond-decorrelation}) are proved
in \cref{app:peyre}. Everything else is proved here from standard facts: the
cyclicity of $\Fp^\times$, the Chinese remainder theorem, Hensel's lemma,
Fourier analysis on finite abelian groups and the Cauchy--Schwarz inequality.
